\section{Durchführung}
\label{sec:Durchführung}

\subsection{Funktionsweise Oszilloskop}
Das Oszilloskop macht zeitlich veränderliche Spannungsverläufe sichtbar. Es kann genau wie ein Multimeter, auch elektrische Gleichspannungen messen.\\
Die Spannungsverläufe werden als Kurve auf einem Bildschirm dargestellt, wobei die X-Achse die Zeit und die Y-Achse die Spannungshöhe abbildet. \\
Man unterscheidet zwischen analoge und digitale Oszilloskope. Ein analoges Oszilloskop nutzt eine Kathodenstrahlröhre, in der ein Elektronenstrahl durch angelegte Spannungen an den X- und Y-Ablenkplatten gezielt gelenkt und auf einem Phosphorschirm sichtbar werden. \\
Die Eingänge des Oszilloskops sind in From von BNC-Buchsen (Bayonet Neill Concelman) vorhanden. Somit wird das Spannungssignal immer relativ zum Erdungspotential gemessen.\\
Da das Oszilloskop nicht immer direkt an das Messobjekt angeschlossen werden kann, benötigt man ein Tastknopf. Der Tastknopf wird dann zwischen den beiden Elementen geschaltet. Er besteht aus einer Messspitze und einer Erdungsklemme. Der typische Innenwiederstand beträgt $10\unit{\mega\ohm}$, damit die Spannungsverläufe bei der Messung nicht verändert werde.\\
Um bei periodischen Signalen ein stehendes Bild zu erhalten, benötigt man einen sogenannten Trigger. Der Trigger legt einen Zeitpunkt fest, wann das Signal auf dem Bildschirm starten soll.\\


\subsection{Aufbau}
Für unsern Versuch bzw. unsere Versuche benötigen wir ein analoges Oszilloskop und eine Leiterplatte, die mit $9\unit{\volt}$ AC Spannung betrieben wird. Auf der Leiterplatte kann man an verschiedenen Stellen Signale mit dem Tastkopf abgreifen. \\
Die Abgriffe 2 und 4 stellen sind Erdungspunkte. Am Abgriff 3 kann man über ein Potentiometer RV1 das Sinussignal für die Amplitude ändern. 


\subsection{Versuchsdurchführung}
Als erstes schalten wir das Oszilloskop ein und schließen die Leiterplatte an die Stromversorgung an.\\
\subsection*{Aufgabe 1: Vermessen einer Rechteckschwingung fester Frequenz}
Der Tastkopf wird geerdet mit der BNC Buchse und die Messspitze wird mit dem 0.2VPP Abgriff verbunden. Der Schalter am Tastkopf steht auf x1. Nun wird das Bild zum Stehen gebracht und die Amplitude Upp, welche von Peak zu Peak geht, wird auf 0.2V eingestellt.\\
\subsection*{Aufgabe 2: Vermessen einer Sinusschwingung fester Frequenz}
Der Schalter des Tastkopfes bleibt weiterhin auf x1 eingestellt. Die Messspitze des Tastkopfes wird mit Abgriff 1 und die Erdung mit Abgriff 2 verbunden. Daraufhin wird das entstandene Bild erneut zum Stehen gebracht, sodass die Amplitude und Periodendauer abgelesen werden kann.\\
\subsection*{Aufgabe 3: Vermessen einer Sinusschwingung variabler Amplitude}
Der Schalter des Tastkopfes bleibt weiterhin auf x1 eingestellt. Die Messspitze wird mit Abgriff 3 und die Erdung bleibt mit Abgriff 2 verbunden. Mit Hilfe des Potentiometers RV! wird die maximale und minimale Amplitude bestimmt. Nun wird der Schalter am Tastkopf auf x10 gestellt und der Vorgang wird wiederholt.\\
\subsection*{Aufgabe 4: Gleichrichter}
Der Schalter am Tastkopf wird wieder auf x1 gestellt. Die Messspitze wird mit Abgriff 5 und die Erdung mit Abgriff 4 verbunden. Es wird auf Einweggleichrichter geschaltet indem der Schalter SW1 in die obere Stellung gebracht wird. Die Einstellungen werden so angepasst und notiert, dass genau drei Halbwellen auf dem Oszilloskop sichtbar sind. Danach wird zunächst auf die positive und dann auf die negative Flanke getriggert.\\
Für den Zweiweggleichrichter wird der Schalter SW1 nun in die untere Stellung gebracht und der gleiche Vorgang, wie für den Einweggleichrichter, wird wiederholt.\\
\subsection*{Aufgabe 5: Gleichstrommessung}
Der Schalter des Tastkopfes bleibt weiterhin mit x1 eingestellt. Die Messspitze wird mit Abgriff 6 und die Erdung bleibt weiterhin mit Abgriff 4 verbunden. Der Schalter SW1 ist in oberer Stellung auf Einweggleichrichter. Der Trigger-Mode wird auf DC geschaltet. Für Ground wird nun der Knopf GD gedrückt und es wird die y-Auslenkung in Bezug aud die Nulllinie gemessen. Die Messspitze des Tastkopfes wird nun mit dem Abgriff 7 und danach mit dem Abgriff 8 verbunden, wo jeweils erneut der gleiche Vorgang wiederholt wird.\\
Der TRigger-Mode wird jetzt aufd AC eingestellt und die Messpitze mit dem Abgriff 6 verbunden. Jetzt wird die Amplitude eines Rippels gemessen. Nun wird auf den zweiweggleichrichter umgeschaltet und der Wert erneut bestimmt.\\
\subsection*{Aufgabe 6: Vermessen eines Signals mit variabler Frequenz}
Der Schalter des Tastkopfes bleibt auf x1 eingestellt. Wir klemmen die Messpitze am Abgriff 9 und die Erdungklemme an Abriff 4 an. Am Potiometer verstellen wir die Frequenz an den Messpunkten 9 und 10. \\